\documentclass[11pt]{article}

\usepackage[margin=1in]{geometry}
\usepackage{amsmath, amssymb}
\usepackage{booktabs}
\usepackage{array}
\usepackage{enumitem}
\usepackage{hyperref}
\usepackage{graphicx}
\usepackage{xcolor}
\usepackage{float}
\usepackage{pgfplots}
\usepackage{pgfplotstable}

\pgfplotsset{compat=1.18}

\hypersetup{
  colorlinks=true,
  linkcolor=blue!60!black,
  urlcolor=blue!60!black,
  citecolor=blue!60!black
}

\newcommand{\coursename}{ECE6323 -- Power System Protection}
\newcommand{\projecttitle}{The Intelligent Merging Unit: Error Correction by Dynamic State Estimation}
\newcommand{\studentname}{Lucas W Peacock}
\newcommand{\partnername}{Leah Kathryn Bryan}
\newcommand{\sectionname}{ECE6323-Q}
\newcommand{\submissiondate}{April 2026}

\title{\projecttitle \\[0.4em] \large Final Report}
\author{\studentname \\ \partnername \\ \coursename \\ \sectionname}
\date{\submissiondate}

\begin{document}
\maketitle
\tableofcontents
\newpage

\section{Introduction and Scope}
The purpose of this project is to estimate the \emph{primary current} of a current transformer (CT)
from the measured burden-side signal available at the merging unit. In the provided data set, the
only directly measured analog quantity is the burden voltage, denoted by $v_{\text{out}}(t)$. The
project objective is to use a model of the CT instrumentation channel together with dynamic state
estimation (DSE) to reconstruct the primary current waveform and other internal states of the
measurement channel.

The instrumentation channel considered in this assignment includes the current transformer, the
interconnecting cable, and the burden/input circuit of the merging unit. The assignment emphasizes
that instrumentation-channel errors can distort measured waveforms and, in severe cases, contribute
to protection-system misoperation. The practical motivation is therefore to create an
``intelligent merging unit'' that can stream corrected estimates of the primary electrical quantity
instead of simply forwarding the raw secondary-side measurement.

This report covers the full implementation requested by the assignment:
\begin{enumerate}[label=\textbf{Task \arabic*:}, leftmargin=1.25in]
  \item Development of a simple user-facing output workflow for inspecting data and results.
  \item Reading COMTRADE files.
  \item Formulating the current instrumentation-channel state model.
  \item Building actual, virtual, derived, and pseudo measurements.
  \item Solving the dynamic state estimation problem.
  \item Processing the provided event files and summarizing the estimated states.
\end{enumerate}

The final implementation is aimed at producing:
\begin{itemize}
  \item the estimated primary current $i_p(t)$,
  \item estimated internal CT flux $\lambda(t)$,
  \item internal channel states such as node voltages and branch currents,
  \item consistency/confidence metrics,
  \item a visual dashboard for the two provided test events.
\end{itemize}

\section{Problem Formulation}
\subsection{Measurement Setting}
The current instrumentation channel is observed at the burden resistor of the merging unit. The
available COMTRADE files contain one analog channel labeled \texttt{Burden\_Voltage}. Therefore,
the estimator does not receive direct measurements of primary current, secondary current, or line
voltage. Instead, those quantities must be reconstructed from the burden voltage and the channel
model.

For the current instrumentation channel, the assignment defines the instrumentation error as
\begin{equation}
\varepsilon = \frac{\left| k\,i_s(t) - i_p(t) \right|}{I_{\text{range}}},
\end{equation}
where $i_s(t)$ is the secondary current, $i_p(t)$ is the primary current, $k$ is the nominal CT
ratio, and $I_{\text{range}}$ is the primary-side measurement range. The project focuses on
reducing this error by estimating the best primary-side state that explains the burden-side
measurement.

\subsection{State Vector}
Following the quadratized model in Appendix A, the estimator uses a 15-state vector:
\begin{equation}
\mathbf{x}(t) =
\begin{bmatrix}
v_1 & v_2 & v_3 & v_4 & e & \lambda & y_1 & y_2 & y_3 & y_4 & i_p & i_m & i_{L1} & i_{L2} & i_{L3}
\end{bmatrix}^{T}.
\end{equation}

The states have the following interpretation:
\begin{itemize}
  \item $v_1$ through $v_4$: internal node voltages of the equivalent circuit,
  \item $e$: transformer-induced voltage,
  \item $\lambda$: CT core magnetic flux,
  \item $i_p$: primary current to be estimated,
  \item $i_m$: magnetizing current,
  \item $i_{L1}$, $i_{L2}$, $i_{L3}$: branch currents in the instrumentation channel,
  \item $y_1$ through $y_4$: auxiliary variables introduced to quadratize the nonlinear
  magnetizing-current relationship.
\end{itemize}

\subsection{Measurement Models}
The assignment specifies 21 equations in the quadratized formulation:
\begin{itemize}
  \item 1 actual measurement,
  \item 10 virtual circuit equations,
  \item 5 derived measurements,
  \item 1 pseudo-measurement,
  \item 4 quadratization relationships.
\end{itemize}

The actual measurement is
\begin{equation}
v_{\text{out}}(t) = v_3(t) - v_4(t).
\end{equation}

The estimator also uses Kirchhoff current law (KCL), Kirchhoff voltage law (KVL), the flux
equation, and the quadratized magnetizing branch relationships as virtual/derived constraints.
Because the project has only one direct sensor waveform, these model equations are essential for
making the state observable.

\subsection{Weighted Least-Squares Dynamic State Estimation}
Let the measurement vector be $\mathbf{z}$ and the nonlinear model be $\mathbf{h}(\mathbf{x})$.
The weighted least-squares objective is
\begin{equation}
J(\mathbf{x}) = \left( \mathbf{z} - \mathbf{h}(\mathbf{x}) \right)^T
\mathbf{W}
\left( \mathbf{z} - \mathbf{h}(\mathbf{x}) \right),
\end{equation}
where $\mathbf{W}$ is a diagonal matrix formed from the inverse variances of the measurement
errors.

The state update is computed using a Gauss--Newton step:
\begin{equation}
\hat{\mathbf{x}}_{k+1}
=
\hat{\mathbf{x}}_{k}
+
\left( \mathbf{H}^{T}\mathbf{W}\mathbf{H} \right)^{-1}
\mathbf{H}^{T}\mathbf{W}
\left( \mathbf{z} - \mathbf{h}(\hat{\mathbf{x}}_{k}) \right),
\end{equation}
where $\mathbf{H}$ is the Jacobian matrix of the measurement equations.

\subsection{Expected Outputs}
The main output of interest is the estimated primary current waveform $i_p(t)$. Additional outputs
include:
\begin{itemize}
  \item estimated CT flux $\lambda(t)$,
  \item estimated internal node voltages and branch currents,
  \item derived burden-current consistency,
  \item confidence/residual metrics,
  \item visual plots for the two test events.
\end{itemize}

\section{Solution Methodology -- Description of the Algorithm}
\subsection{Task 0: User-Facing Output and Visualization}
To satisfy the requirement for reporting simulated data, intermediate values, and relay-relevant
results, the solution produces:
\begin{itemize}
  \item \texttt{results.csv} containing per-sample estimated quantities,
  \item \texttt{residuals.csv} containing per-sample residuals for all measurement equations,
  \item \texttt{summary.json} containing aggregate statistics,
  \item an HTML dashboard with tabs for Event 01 and Event 02.
\end{itemize}
The dashboard provides a compact visual summary of the burden voltage, estimated primary current,
estimated flux, confidence, and normalized residual behavior for each event.

\subsection{Task 1: COMTRADE Reader}
The implementation reads the COMTRADE configuration (\texttt{.cfg}) file to determine:
\begin{itemize}
  \item number of analog channels,
  \item channel names and units,
  \item scale factors $a$ and offsets $b$,
  \item sample rate and file type.
\end{itemize}

The ASCII \texttt{.dat} file is then parsed line by line. Each raw analog sample is converted to a
physical value using
\begin{equation}
x_{\text{physical}} = a x_{\text{raw}} + b.
\end{equation}

For both provided events, the parser confirmed that the data set contains only one analog channel:
\texttt{Burden\_Voltage}.

\subsection{Task 2: Mathematical Model of the Instrumentation Channel}
The channel is modeled with the assignment parameters:
\begin{itemize}
  \item CT turns ratio $n = 400$,
  \item inductances $L_1$, $L_2$, $L_3$ and mutual inductance $M_{23}$,
  \item burden conductance $g_b = 1/R_b$ with $R_b = 0.1 \, \Omega$,
  \item magnetizing branch parameters $(i_0, \lambda_0, L_0)$,
  \item cable and winding resistances $r_1$, $r_2$, $r_3$.
\end{itemize}

The model uses the equivalent-circuit relationships given in Appendix A. The nonlinear
magnetizing-current relation is quadratized using the auxiliary variables $y_1$ through $y_4$ so
that the final estimator operates on a quadratic-algebraic measurement model.

\subsection{Task 3: Measurement Construction}
The algorithm builds a 21-element measurement vector at each sample. It contains:
\begin{enumerate}[label=\alph*)]
  \item the actual burden-voltage measurement,
  \item zero-valued virtual measurements enforcing KCL/KVL and flux equations,
  \item derived burden-current measurements implied by the burden voltage,
  \item a pseudo-measurement that anchors the grounded node.
\end{enumerate}

Because only one direct waveform is available, the virtual and derived equations are important.
They do not add new sensor information, but they encode the physical laws that the estimated state
must obey.

\subsection{Task 4: Dynamic State Estimation Algorithm}
The estimator proceeds sample by sample. At each time step:
\begin{enumerate}
  \item the burden-voltage sample is read from the COMTRADE data,
  \item a measurement vector is formed,
  \item an initial state guess is constructed from the previous state and the burden current
  implied by $v_{\text{out}}$,
  \item the nonlinear measurement model is solved with a damped Gauss--Newton iteration,
  \item a chi-square style consistency metric is computed from the weighted residuals.
\end{enumerate}

To represent the dynamic terms, derivatives are approximated using the previous solved state and
the current sample interval:
\begin{equation}
\frac{d x(t)}{dt} \approx \frac{x_k - x_{k-1}}{\Delta t}.
\end{equation}

This sequential update makes the estimator a practical dynamic solver while still keeping the
implementation computationally light enough for the project data set.

\subsection{Task 5: Event Processing Workflow}
The two provided COMTRADE events were processed independently. The algorithm reads one event at a
time, evaluates every sample sequentially, records the estimated state, and then writes the
results to output files for later analysis and visualization.

\subsection{Task 6: Summary Variables Reported}
The reported outputs include:
\begin{itemize}
  \item estimated primary current $i_p(t)$,
  \item estimated magnetic flux $\lambda(t)$,
  \item estimated induced voltage $e(t)$,
  \item sample-wise consistency/confidence values,
  \item maximum normalized residuals,
  \item event summary statistics.
\end{itemize}

\section{Presentation of Results}
\subsection{Input Events}
The assignment provides two fault/disturbance events with different degrees of CT saturation. Both
events contain 4700 samples at approximately 4.8 kHz. Event 02 is visibly more severe than Event
01 based on the burden-voltage swing and the resulting internal-state magnitudes.

\subsection{Event Summary}
Table~\ref{tab:summary} summarizes the principal numerical outputs from the estimator.

\begin{table}[H]
  \centering
  \caption{Summary of estimated event results}
  \label{tab:summary}
  \begin{tabular}{>{\raggedright\arraybackslash}p{2.6in}cc}
    \toprule
    Metric & Event 01 & Event 02 \\
    \midrule
    Number of samples & 4700 & 4700 \\
    Converged samples & 4700 & 4700 \\
    Mean confidence & 1.0000 & 1.0000 \\
    Minimum confidence & 1.0000 & 1.0000 \\
    Maximum estimated primary current (A) & 4190.121 & 19052.104 \\
    Minimum estimated primary current (A) & -5857.761 & -27757.416 \\
    Maximum estimated flux (Wb) & 0.008401 & 0.087216 \\
    Maximum absolute normalized residual & $1.55 \times 10^{-5}$ & $1.75 \times 10^{-5}$ \\
    \bottomrule
  \end{tabular}
\end{table}

\subsection{Discussion of Results}
Several observations are immediate:
\begin{enumerate}
  \item Event 02 produces a much larger estimated primary-current excursion than Event 01.
  \item Event 02 also produces significantly higher estimated CT flux, indicating a more severe
  saturation case.
  \item The burden-voltage waveform alone is sufficient, in this model-based framework, to drive a
  full state estimate of the instrumentation channel.
  \item The confidence metric is very high after enforcing the model constraints consistently.
\end{enumerate}

The most important engineering result is that Event 02 stresses the instrumentation channel much
more strongly than Event 01. This agrees with the assignment statement that the two files were
constructed to represent different degrees of CT saturation.

\subsection{Event Waveforms and Chi-Square Trends}
Figures~\ref{fig:burden-voltage-events} and~\ref{fig:chi-square-events} provide direct visual
evidence of the difference in severity between the two events. The burden-voltage plot shows that
Event 02 has a substantially larger excursion than Event 01, which is consistent with a more
heavily stressed instrumentation channel. The chi-square plot shows the numerical consistency of
the solved estimate at each time sample.

\begin{figure}[H]
  \centering
  \begin{tikzpicture}
    \begin{axis}[
      width=0.9\textwidth,
      height=2.8in,
      grid=both,
      grid style={line width=0.1pt, draw=gray!20},
      major grid style={line width=0.2pt, draw=gray!35},
      xlabel={Time (s)},
      ylabel={Burden Voltage (V)},
      title={Measured Burden Voltage for Event 01 and Event 02},
      legend style={at={(0.02,0.98)}, anchor=north west, draw=none, fill=white},
      yticklabel style={/pgf/number format/fixed},
      tick label style={font=\small},
      label style={font=\small},
      title style={font=\normalsize}
    ]
      \addplot[
        color=blue!70!black,
        line width=0.8pt
      ] table[
        col sep=comma,
        header=true,
        x index=1,
        y index=2
      ]{outputs/event01/results.csv};
      \addlegendentry{Event 01}

      \addplot[
        color=red!75!black,
        line width=0.8pt
      ] table[
        col sep=comma,
        header=true,
        x index=1,
        y index=2
      ]{outputs/event02/results.csv};
      \addlegendentry{Event 02}
    \end{axis}
  \end{tikzpicture}
  \caption{Measured burden-voltage waveforms from the two provided COMTRADE events. Event 02 has
  a much larger magnitude swing than Event 01, indicating a more severe instrumentation-channel
  disturbance.}
  \label{fig:burden-voltage-events}
\end{figure}

\begin{figure}[H]
  \centering
  \begin{tikzpicture}
    \begin{axis}[
      width=0.9\textwidth,
      height=2.8in,
      grid=both,
      grid style={line width=0.1pt, draw=gray!20},
      major grid style={line width=0.2pt, draw=gray!35},
      xlabel={Time (s)},
      ylabel={Chi-square value},
      title={Chi-Square Consistency of the Dynamic State Estimation},
      legend style={at={(0.02,0.98)}, anchor=north west, draw=none, fill=white},
      ymode=log,
      log basis y={10},
      tick label style={font=\small},
      label style={font=\small},
      title style={font=\normalsize}
    ]
      \addplot[
        color=blue!70!black,
        line width=0.8pt
      ] table[
        col sep=comma,
        header=true,
        x index=1,
        y index=8
      ]{outputs/event01/results.csv};
      \addlegendentry{Event 01}

      \addplot[
        color=red!75!black,
        line width=0.8pt
      ] table[
        col sep=comma,
        header=true,
        x index=1,
        y index=8
      ]{outputs/event02/results.csv};
      \addlegendentry{Event 02}
    \end{axis}
  \end{tikzpicture}
  \caption{Per-sample chi-square values of the estimation. In this implementation, the chi-square
  values remain very small after the measurement-model sign correction, indicating that the solved
  state is highly consistent with the selected model equations.}
  \label{fig:chi-square-events}
\end{figure}

\subsection{Interpretation of the Confidence Metric}
The confidence values computed in this implementation should be interpreted carefully. The data set
contains one direct physical waveform and many model-based virtual or derived equations. Therefore,
the confidence metric is best interpreted as a \emph{model-consistency indicator} rather than a
stand-alone proof of physical ground truth. In other words, a high confidence means that the state
estimate satisfies the selected measurement model very well; it does not automatically mean that
the true field current is independently verified by multiple sensors.

\subsection{Visualization Deliverables}
The implementation produces a tabbed HTML dashboard with separate tabs for Event 01 and Event 02.
Each tab contains:
\begin{itemize}
  \item the burden voltage waveform,
  \item estimated primary current,
  \item estimated core flux,
  \item confidence over time,
  \item maximum normalized residual over time,
  \item summary cards for the event statistics.
\end{itemize}

This dashboard is useful for quickly comparing the two event severities and for validating that the
estimator outputs remain physically and numerically consistent over the full sample window.

\section{Conclusions and Observations}
This project implemented a dynamic state estimator for CT instrumentation-channel error correction
using only the burden-side voltage waveform from COMTRADE data and the provided channel model. The
solution successfully completed the required workflow:
\begin{itemize}
  \item COMTRADE data were parsed correctly.
  \item The current instrumentation channel was modeled with the provided equivalent circuit.
  \item Actual, virtual, derived, and pseudo measurements were constructed.
  \item A Gauss--Newton weighted least-squares dynamic estimator was implemented.
  \item Both provided events were processed and summarized.
  \item A visual dashboard was generated for event-level interpretation.
\end{itemize}

The recovered results show that Event 02 is the more severe case, producing much higher estimated
primary current and much larger core flux than Event 01. This is consistent with the assignment's
description of varying CT saturation severity across the events.

The added waveform and chi-square figures support the same conclusion visually. The burden-voltage
plot shows a much stronger excursion for Event 02, while the chi-square plot shows that both event
solutions remain numerically consistent with the final estimator model over the full duration of
the record.

The main limitation of the problem is that only one direct analog waveform is available. Because of
this, the quality of the final estimate depends strongly on the correctness of the instrumentation
channel model and the selected virtual/derived measurement set. Nonetheless, this is also the main
point of the assignment: to show that a merging unit can use model-based DSE to reconstruct a more
useful primary-side estimate from limited secondary-side data.

Possible future improvements include:
\begin{itemize}
  \item adding independent validation data, if available,
  \item comparing multiple derivative-discretization schemes,
  \item extending the same methodology to a voltage instrumentation channel,
  \item integrating the estimator more directly with a protection or relay decision block.
\end{itemize}

\begin{thebibliography}{9}
\bibitem{meliopoulos2018}
Sakis Meliopoulos, George Cokkinides, Jiahao Xie, and Yuan Kong,
``Instrumentation Error Correction within Merging Units,''
\emph{Proceedings of the 2018 Georgia Tech Fault and Disturbance Analysis Conference},
Atlanta, Georgia, April 30--May 1, 2018.

\bibitem{meliopoulos2019}
Sakis Meliopoulos, Jiahao Xie, and George Cokkinides,
``Impact of Error Correction on Protection Performance during Geomagnetic Disturbances,''
\emph{Proceedings of the 2019 Georgia Tech Fault and Disturbance Analysis Conference},
Atlanta, Georgia, April 29--30, 2019.

\bibitem{kong2021}
Yuan Kong and A. P. Sakis Meliopoulos,
``Intelligent Merging Unit: DSE-based Online Instrumentation Channel Error Correction,''
\emph{Electric Power Systems Research},
Volume 200, November 2021, 107468.

\bibitem{obikwelu2021}
Emeka Obikwelu and A. P. Sakis Meliopoulos,
``VT Instrumentation Channel Error Correction using Dynamic State Estimation,''
\emph{Proceedings of the 52nd North American Power Symposium},
Phoenix, Arizona, April 11--13, 2021.
\end{thebibliography}

\end{document}
